\subsection{Definição de Estados e Variáveis de Controle}
\label{sec:var_estado_controle}

Para este trabalho, um conjunto de variáveis de estado que descrevem a dinâmica orbital, a atitude e os elos do MR foram estabelecidos, e são:

\begin{equation}
    \label{eq:variaveis_estado}
    \mathbf{x} =
    \{ \vec r,\; \vec v,\; q,\; \vec \omega,\; \theta,\; \dot\theta \},
\end{equation}

onde $\vec r$ e $\vec v$ representam a posição e a velocidade relativas, expressas no referencial LVLH;
$q$ corresponde ao quaternion unitário que descreve a atitude da espaçonave visitante em relação ao referencial inercial;
$\vec \omega$ é o vetor de velocidade angular da base da espaçonave;
$\theta$ indica os ângulos articulares do MR;
e $\dot \theta$ representa as respectivas velocidades articulares.
As variáveis de controle são representadas por forças e torques:

\begin{equation}
    \label{eq:variaveis_controle}
    \mathbf{u} =
    \{ \vec f,\; \vec \tau \},
\end{equation}

onde $\vec f$ são as forças de propulsão que atuam no controle orbital, tanto absoluto quanto relativo, e $\vec \tau$ são os torques aplicados para o controle de atitude da espaçonave e para o acionamento das juntas do manipulador.

O modelo dinâmico considerado integra diferentes subsistemas, como:
(i) as equações de Hill/Clohessy--Wiltshire (HCW) linearizadas para a descrição do movimento relativo,
(ii) as equações de atitude baseadas em formulações de Euler e quaternions, e
(iii) as equações dinâmicas do manipulador robótico acoplado à base da espaçonave.
